\documentclass[10pt]{article}
\usepackage[T1]{fontenc}
\usepackage[utf8]{inputenc}
\usepackage[french]{babel}
\usepackage{amsmath,amssymb,mathrsfs,mathpazo}
\usepackage{geometry}
\geometry{a4paper,margin=2cm}
\usepackage{graphicx,array,tabularx,tabularray,colortbl,multirow,multicol,booktabs,enumitem,pifont,ulem,fancybox,diagbox,cancel}
\usepackage{pstricks,pstricks-add,pst-plot,pst-tree,pst-node,pst-text,pst-func,pst-math,pst-bspline,pst-3dplot,pst-eucl,pst-solides3d}
\let\pscalculate\undefined
\usepackage{pst-bezier}
\usepackage{tikz,pgfplots,tkz-tab}
\usetikzlibrary{arrows,arrows.meta,patterns,calc,positioning,trees,decorations.markings,shapes.geometric,shapes.misc,shapes.arrows,intersections,decorations.pathmorphing,decorations.pathreplacing,backgrounds,shadows,babel}
\pgfplotsset{compat=1.18}
\usepgfplotslibrary{fillbetween}
\usepackage[np]{numprint}
\usepackage{listings,makecell,pas-tableur,twoopt}
\let\barre\undefined
\usepackage[pstricks]{tablvar}
\definecolor{darkgreen}{RGB}{0,100,0}
\providecommand{\red}{\color{red}}
\providecommand{\blue}{\color{blue}}
\providecommand{\green}{\color{green}}
\providecommand{\black}{\color{black}}
\providecommand{\white}{\color{white}}
\providecommand{\gray}{\color{gray}}
\providecommand{\cyan}{\color{cyan}}
\providecommand{\magenta}{\color{magenta}}
\providecommand{\yellow}{\color{yellow}}
\providecommand{\violet}{\color{violet}}
\providecommand{\orange}{\color{orange}}
\providecommand{\euro}{\texteuro}
\providecommand{\EUR}{\texteuro}
\providecommand{\ieme}{\textsuperscript{e}}
\providecommand{\degres}{\ensuremath{^\circ}}
\providecommand{\degree}{\ensuremath{^\circ}}
\providecommand{\vv}[1]{\overrightarrow{#1}}
\providecommand{\Coord}[1]{#1}
\providecommand{\decosix}{}
\providecommand{\starredbullet}{\ensuremath{\bullet}}
\pagestyle{empty}
\setlength{\parindent}{0pt}
\usepackage{adjustbox}\DeclareUnicodeCharacter{2236}{:}\DeclareUnicodeCharacter{2190}{\ensuremath{\leftarrow}}
\providecommand{\R}{\mathbb{R}}
\providecommand{\N}{\mathbb{N}}
\providecommand{\Z}{\mathbb{Z}}
\providecommand{\Q}{\mathbb{Q}}
\providecommand{\D}{\mathbb{D}}
\providecommand{\C}{\mathbb{C}}
\providecommand{\e}{{\mathrm{e}}}
\providecommand{\vect}[1]{\overrightarrow{#1}}
\providecommand{\vectt}[1]{\overrightarrow{\mathrm{#1}}}
\providecommand{\vvt}[1]{\overrightarrow{\mathrm{#1}}}
\providecommand{\np}[2][]{{#2~#1}}
\providecommand{\vv}[1]{{\overrightarrow{#1}}}
\renewcommand{\d}{{\mathrm{d}}}
\providecommand{\degres}{{^\circ}}
\providecommand{\degre}{{^\circ}}
\providecommand{\blue}{\color{blue}}
\providecommand{\red}{\color{red}}
\providecommand{\enskip}{\quad}
\providecommand{\llbracket}{\lbrack\!\lbrack}
\providecommand{\rrbracket}{\rbrack\!\rbrack}
\providecommand{\euro}{{\text{€}}}
\providecommand{\no}{{\text{n°}}}
\providecommand{\interval}[2]{{[#1;#2]}}
\providecommand{\intervFO}[2]{{[#1;#2[}}
\providecommand{\dotfill}{\ldots}

\providecommand{\ds}{\displaystyle}
\providecommand{\textding}[1]{#1}
\providecommand{\Oij}{\(\mathrm{O};\vec{\imath},\vec{\jmath}\)}
\providecommand{\Oijk}{\(\mathrm{O};\vec{\imath},\vec{\jmath},\vec{k}\)}
\providecommand{\Ouv}{\(\mathrm{O};\vec{u},\vec{v}\)}
\providecommand{\intervFF}[2]{\ensuremath{[#1;#2]}}
\renewcommand{\euro}{\ifmmode\text{\texteuro}\else\texteuro\fi}
\DeclareUnicodeCharacter{00B0}{\ifmmode^{\circ}\else\textdegree\fi}
\DeclareUnicodeCharacter{20AC}{\ifmmode\text{\texteuro}\else\texteuro\fi}
% Énoncés : Éducation nationale. Transcription et corrigés : APMEP, Denis Vergès et contributeurs. Découpage et classement : CoopMaths, Sébastien Lozano, Éric Elter et contributeurs. Adaptation de la présentation : Maths Explorer.
% CC BY-SA : https://coopmaths.fr/www/about/
% https://coopmaths.fr/alea/static/bac/2024/tex/bac_2024_06_sujet1_metropole_devoile_4_cor.tex
\begin{document}
\noindent\textit{APMEP / CoopMaths. CC BY-SA. Présentation adaptée par Maths Explorer.}\par\medskip
\medskip






	\begin{enumerate}
		\item On considère ci-dessous le tableau de variations d'une fonction $f$ définie sur $\R\setminus \{- 2\}$.

\begin{center}
{\renewcommand{\arraystretch}{1.1}
\psset{nodesep=3pt,arrowsize=2pt 3}  
\def\esp{\hspace*{1.5cm}}
\def\hauteur{0pt}
$\begin{array}{|c| c c c r c l c c c |}
\hline
 x & -\infty & \esp &  & -2 & & \esp & 1 & \esp & +\infty \\
\hline
  &   \Rnode{max1}{5} & &   & \vline\;\vline &  & & \Rnode{max2}{3} & & \\
f  & &  &  & \vline\;\vline & & & & & \rule{0pt}{\hauteur}\\
 & & & \Rnode{min1}{-\infty}  & \vline\;\vline & \Rnode{min2}{-\infty} & &  & & \Rnode{min3}{1} \rule{0pt}{\hauteur} 
\ncline{->}{max1}{min1} 
\ncline{->}{min2}{max2}
\ncline{->}{max2}{min3} \\
\hline
\end{array}$
}
\end{center}
	\begin{enumerate}
		\item \textbf{Affirmation 1 :}
\og La droite d'équation $y = -2$ est asymptote horizontale à la courbe $\mathcal{C}_f$ de la fonction $f$. \fg{} 

La droite d'équation $y = -2$ est asymptote horizontale à la courbe $\mathcal{C}_f$ de la fonction $f$ si et seulement si $\ds\lim_{x\to -\infty} f(x)=-2$ ou $\ds\lim_{x\to +\infty} f(x)=-2$.



Or, d'après le tableau de variations de $f$, $\ds\lim_{x\to -\infty} f(x)=5$ et $\ds\lim_{x\to +\infty} f(x)=1$.

\hfill\textbf{Affirmation 1 fausse}

		\item \textbf{Affirmation 2:}
\og  $\displaystyle\lim_{x \to - \infty} \dfrac{2}{f(x) - 5} = + \infty$. \fg{}

D'après le tableau de variations de $f$, celle-ci est décroissante sur l'intervalle $]- \infty~;~-2[$, donc sur cet intervalle $f(x) < 5$ et donc $f(x)-5<0$.

Comme $\displaystyle\lim_{x \to - \infty} f(x) = 5$, on a donc $\displaystyle\lim_{x \to - \infty} f(x) - 5 = 0^{-}$ et donc\\
 $\displaystyle\lim_{x \to - \infty} \dfrac{1}{f(x) - 5} = - \infty$ et enfin  $\displaystyle\lim_{x \to - \infty} \dfrac{2}{f(x) - 5} = - \infty$.
 
 \hfill\textbf{Affirmation 2 fausse}
	\end{enumerate}
	
\item On considère la fonction $g$ définie sur $\R$ par $g(x) = x\e^{-x}$.
	\begin{enumerate}
	
		\item \textbf{Affirmation 3 :} 
\og  Le point A$\left(2~;~\dfrac{2}{\e^2}\right)$ est l'unique point d'inflexion de la courbe $\mathcal{C}_g$ de la fonction~$g$. \fg{} 

La fonction $g$ produit de fonctions dérivables sur $\R$ est dérivable et sur cet intervalle :

$g'(x) = \e^{-x} - x\e^{-x} = \e^{-x}(1 - x)$, qui est elle-même dérivable sur $\R$ et 

$g''(x) = -\e^{-x} - (1 - x)\e^{-x} = \e^{-x}(x - 2)$.

On a $g''(x) = 0 \iff \e^{-x}(x - 2) \iff x - 2 = 0$, puisque $\e^{-x} \ne 0$, quel que soit $x \in \R$, donc la dérivée seconde ne s'annule qu'en $x = 2$.

Comme $g(2) = 2\e^{-2} = \dfrac{2}{\e^2}$ et comme $g''(x)$ a le signe de $(x - 2)$ et donc change de signe (de moins à plus) en 2, le point A$\left(2~;~\dfrac{2}{\e^2}\right)$ est l'unique point d'inflexion de la courbe $\mathcal{C}_g$.

 \hfill\textbf{Affirmation 3 vraie}
 
		\item \textbf{Affirmation 4 :} 
\og  Pour tout nombre réel $x$ appartenant à $] - \infty~;~2[$, on a $g(x) \leqslant x$. \fg{} 

$g(x) \leqslant x \iff x\e^{-x} \leqslant x \iff x \leqslant x\e^{x} \iff 0 \leqslant x\e^{x}-x \iff x\left (\e^{x}-1\right )\geqslant 0$

On établit un tableau de signes:

\begin{center}
{
\def\esp{\hspace*{3cm}}
$\begin{array}{|c | *{5}{c} |} 
\hline
x  & -\infty & \esp & 0 & \esp  & 2 \\
\hline
x &  & \pmb{-} &  \vline\hspace{-2.7pt}{0} & \pmb{+} &   \\
\hline
\e^{x}-1 &  & \pmb{-} &  \vline\hspace{-2.7pt}{0} & \pmb{+} &   \\
\hline
x\left (\e^{x}-1\right ) &  & \pmb{+} &  \vline\hspace{-2.7pt}{0} & \pmb{+} &   \\
\hline
\end{array}$
}
\end{center}

Donc, pour tout $x$ de $] - \infty~;~2[$, on a $x\left (\e^{x}-1\right ) \geqslant 0$ donc $g(x) \leqslant x$.








 
  \hfill\textbf{Affirmation 4 vraie}
\end{enumerate}
	
\item \textbf{Affirmation 5 :} 
\og  L'équation $x\, \ln (x) = 1$ admet exactement deux solutions sur l'intervalle $]0~;~ +\infty[$. \fg{}

Soit $f$ la fonction définie sur $\left ]0\;;\;+\infty\strut\right [$ par $f(x)=x\,\ln(x)-1$.

$x\, \ln (x) = 1 \iff x\, \ln (x) - 1 = 0 \iff f(x)=0$

\begin{list}{\textbullet}{}
\item La fonction $f$ est dérivable et $f'(x)=1\times \ln(x) + x\times \dfrac{1}{x}-0=\ln(x)+1$.
\item $f'(x)>0 \iff \ln(x) +1>0 \iff \ln(x)>-1 \iff x>\e^{-1}$
\item $\ds\lim_{x\to 0 \atop x>0} x\,\ln(x)=0$ donc $\ds\lim_{x\to 0 \atop x>0} f(x)=-1$
\item $\ds\lim_{x\to +\infty} x =+\infty$ et $\ds\lim_{x\to +\infty} \ln(x)=+\infty$ donc $\ds\lim_{x\to +\infty} f(x)=+\infty$
\item  $f(\e^{-1})=\e^{-1}\,\ln(\e^{-1})-1 = -\e^{-1}-1<0$
\end{list}

D'où le tableau de variation de $f$:

\begin{center}
{\renewcommand{\arraystretch}{1.1}
\psset{nodesep=3pt,arrowsize=2pt 3}  
\def\esp{\hspace*{2cm}}
\def\hauteur{0pt}
$\begin{array}{|c| l  *4{c}|}
\hline
 x & 0 & \esp & \e^{-1} & \esp & +\infty \\
 \hline
f'(x) & \vline\;\vline &  \pmb{-} & \vline\hspace{-2.7pt}0 & \pmb{+} & \\  
\hline
  & \vline\;\vline~\Rnode{max1}{-1}  &  &  &  & \Rnode{max2}{+\infty} \\
f(x) &\vline\;\vline &  & & &  \rule{0pt}{\hauteur} \\
 & \vline\;\vline & &   \Rnode{min}{-\e^{-1}-1} & & \rule{0pt}{\hauteur}
\ncline{->}{max1}{min} \ncline{->}{min}{max2}\\
\hline
\end{array}$
}
\end{center}	

Sur l'intervalle $\left ]0\;;\;\e^{-1}\strut\right [$, la fonction $f$ est négative, donc l'équation $f(x)=0$ n'a pas de solution.

Sur l'intervalle $\left ]\e^{-1}\;;\; +\infty\strut\right [$, la fonction $f$, continue et strictement croissante, passe du négatif au positif: d'après le corollaire du théorème des valeurs intermédiaires, l'équation $f(x)=0$ admet une solution unique.

\hfill\textbf{Affirmation 5 fausse}
\end{enumerate}

\end{document}
